\addcontentsline{toc}{chapter}{Bibliografía}

\begin{thebibliography}{99}
\raggedright

\bibitem{hqcspec}
C.~Aguilar-Melchor, N.~Aragon, S.~Bettaieb, L.~Bidoux, O.~Blazy,
et al. \emph{Hamming Quasi-Cyclic (HQC): specification}, submission
to the NIST post-quantum cryptography standardization process,
round~4, 2023.

\bibitem{nist8545}
G.~Alagic et al. \emph{Status report on the fourth round of the NIST
post-quantum cryptography standardization process}, NIST Internal
Report 8545, March 2025.

\bibitem{apon2026}
D.~Apon. An algebraic-geometry lower bound against the ePrint
2026/1747 McEliece key-recovery attack, Cryptology ePrint Archive,
Report 2026/1810, 2026.

\bibitem{bardet2026}
M.~Bardet, A.~Lemoine, J.-P.~Tillich. An attack on the CFS scheme
and on TII McEliece challenges. In \emph{Advances in Cryptology ---
CRYPTO 2026}, Springer, 2026.

\bibitem{bjmm2012}
A.~Becker, A.~Joux, A.~May, A.~Meurer. Decoding random binary linear
codes in $2^{n/20}$: How $1+1=0$ improves information set decoding.
In \emph{Advances in Cryptology --- EUROCRYPT 2012}, LNCS 7237,
pp.~520--536, Springer, 2012.

\bibitem{bcmv1978}
E.~R.~Berlekamp, R.~J.~McEliece, H.~C.~A.~van Tilborg. On the
inherent intractability of certain coding problems. \emph{IEEE
Transactions on Information Theory} 24(3):384--386, 1978.

\bibitem{bernstein2010}
D.~J.~Bernstein. Grover vs.\ McEliece. In \emph{Post-Quantum
Cryptography --- PQCrypto 2010}, LNCS 6061, pp.~73--80, Springer,
2010.

\bibitem{ntrufault}
D.~J.~Bernstein. A one-time single-bit fault leaks all previous
NTRU-HRSS session keys to a chosen-ciphertext attack, Cryptology
ePrint Archive, Report 2022/1125, 2022.

\bibitem{bernstein2026}
D.~J.~Bernstein et al. Posturas y respuestas del equipo Classic
McEliece, lista de correo del \emph{pqc-forum} del NIST, agosto--octubre
de 2026, \url{https://groups.google.com/a/list.nist.gov/g/pqc-forum}.

\bibitem{cmspec}
D.~J.~Bernstein, T.~Chou, T.~Lange, et al. \emph{Classic McEliece:
conservative code-based cryptography} (specification), 23 de octubre
de 2022, \url{https://classic.mceliece.org/mceliece-spec-20221023.pdf}.

\bibitem{blp2011}
D.~J.~Bernstein, T.~Lange, C.~Peters. Smaller decoding exponents:
ball-collision decoding. In \emph{Advances in Cryptology ---
CRYPTO 2011}, LNCS 6841, pp.~743--760, Springer, 2011.

\bibitem{bp2018}
D.~J.~Bernstein, E.~Persichetti. Towards KEM unification, Cryptology
ePrint Archive, Report 2018/526, 2018.

\bibitem{berson1997}
T.~A.~Berson. Failure of the McEliece public-key cryptosystem under
message-resend and related-message attack. In \emph{Advances in
Cryptology --- CRYPTO~'97}, LNCS 1294, pp.~213--220, Springer, 1997.

\bibitem{bindel2019}
N.~Bindel, M.~Hamburg, K.~H\"ovelmanns, A.~H\"ulsing, E.~Persichetti.
Tighter proofs of CCA security in the quantum random oracle model.
In \emph{Theory of Cryptography --- TCC 2019}, LNCS 11892,
pp.~61--90, Springer, 2019.

\bibitem{bothmay2018}
L.~Both, A.~May. Decoding linear codes with high error rate and its
impact for LPN security. In \emph{Post-Quantum Cryptography ---
PQCrypto 2018}, LNCS 10786, pp.~25--46, Springer, 2018.

\bibitem{briaud2026}
P.~Briaud, A.~Lemoine, H.~Randriambololona, J.-P.~Tillich. A
heuristic subexponential attack on the McEliece cryptosystem. In
\emph{Advances in Cryptology --- ASIACRYPT 2026}, Springer, 2026.

\bibitem{bsi2026}
Bundesamt f\"ur Sicherheit in der Informationstechnik. \emph{Notes on
recent developments concerning Classic McEliece}, 29 de septiembre
de 2026, \url{https://www.bsi.bund.de}.

\bibitem{canteautsendrier1998}
A.~Canteaut, N.~Sendrier. Cryptanalysis of the original McEliece
cryptosystem. In \emph{Advances in Cryptology --- ASIACRYPT~'98},
LNCS 1514, pp.~187--199, Springer, 1998.

\bibitem{carrier2022}
K.~Carrier, T.~Debris-Alazard, C.~Meyer-Hilfiger, J.-P.~Tillich.
Statistical decoding 2.0: reducing decoding to LPN,
arXiv:2208.02201, 2022.

\bibitem{cm529}
Classic McEliece Team. \emph{Notes on a recent claim that a
\texttt{mceliece348864} distinguisher uses only $2^{529}$
operations}, 17 de abril de 2025,
\url{https://classic.mceliece.org/mceliece-529-20250417.pdf}.

\bibitem{cm610}
Classic McEliece Team. \emph{Notes on recent speculation that a
\texttt{mceliece348864} key-recovery attack uses only $2^{610}$
operations}, 23 de junio de 2026,
\url{https://classic.mceliece.org/mceliece-610-20260623.pdf}.

\bibitem{cm2025yodo}
Classic McEliece Team. \emph{Statement regarding the paper ``You Only
Decapsulate Once''}, 20 de noviembre de 2025,
\url{https://classic.mceliece.org/mceliece-20251120.txt}.

\bibitem{cmimpl}
Classic McEliece Team. \emph{Classic McEliece: conservative
code-based cryptography: guide for implementors}, 23 de octubre de
2022,
\url{https://classic.mceliece.org/mceliece-impl-20221023.pdf}.

\bibitem{cmiso}
Classic McEliece Team. P\'agina de estandarizaci\'on ISO,
\url{https://classic.mceliece.org/iso.html}, consultada en octubre de
2026.

\bibitem{cmmods3}
Classic McEliece Team. \emph{Classic McEliece: conservative
code-based cryptography: description of modifications from round 3
to round 4}, 23 de octubre de 2022,
\url{https://classic.mceliece.org/mceliece-mods3-20221023.pdf}.

\bibitem{cmnist}
Classic McEliece Team. P\'agina de estandarizaci\'on NIST,
\url{https://classic.mceliece.org/nist.html}, consultada en octubre de
2026.

\bibitem{cmpc}
Classic McEliece Team. \emph{What plaintext confirmation means},
23 de octubre de 2022,
\url{https://classic.mceliece.org/mceliece-pc-20221023.pdf}.

\bibitem{cmrationale}
Classic McEliece Team. \emph{Classic McEliece: conservative
code-based cryptography: design rationale}, 23 de octubre de 2022,
\url{https://classic.mceliece.org/mceliece-rationale-20221023.pdf}.

\bibitem{cmsecurity}
Classic McEliece Team. \emph{Classic McEliece: conservative
code-based cryptography: guide for security reviewers}, 23 de
octubre de 2022,
\url{https://classic.mceliece.org/mceliece-security-20221023.pdf}.

\bibitem{couvreur2017}
A.~Couvreur, A.~Otmani, J.-P.~Tillich. Polynomial time attack on
wild McEliece over quadratic extensions. \emph{IEEE Transactions on
Information Theory} 63(1):404--427, 2017.

\bibitem{dumer1991}
I.~Dumer. On minimum distance decoding of linear codes. En
\emph{Proc.\ 5th Joint Soviet-Swedish International Workshop on
Information Theory}, pp.~50--52, 1991.

\bibitem{esserbellini2022}
A.~Esser, E.~Bellini. Syndrome decoding estimator. In
\emph{Public-Key Cryptography --- PKC 2022}, LNCS 13177,
pp.~112--141, Springer, 2022.

\bibitem{faugere2013}
J.-C.~Faug\`ere, V.~Gauthier-Uma\~na, A.~Otmani, L.~Perret,
J.-P.~Tillich. A distinguisher for high-rate McEliece cryptosystems.
\emph{IEEE Transactions on Information Theory} 59(10):6830--6844,
2013.

\bibitem{fips203}
National Institute of Standards and Technology. \emph{Module-lattice-based
key-encapsulation mechanism standard}, FIPS 203, agosto de 2024.

\bibitem{fips204}
National Institute of Standards and Technology. \emph{Module-lattice-based
digital signature standard}, FIPS 204, agosto de 2024.

\bibitem{fips205}
National Institute of Standards and Technology. \emph{Stateless
hash-based digital signature standard}, FIPS 205, agosto de 2024.

\bibitem{foj1999}
E.~Fujisaki, T.~Okamoto. Secure integration of asymmetric and
symmetric encryption schemes. In \emph{Advances in Cryptology ---
CRYPTO~'99}, LNCS 1666, pp.~537--554, Springer, 1999.

\bibitem{ghoshal2026}
A.~Ghoshal, Y.~Ishai, A.~Jain, N.~Sun. Quasipolynomial cryptanalysis
of the McEliece cryptosystem (or: PIR meets McEliece), Cryptology
ePrint Archive, Report 2026/1630, 2026.

\bibitem{gidney2021}
C.~Gidney, M.~Eker\r{a}. How to factor 2048 bit RSA integers in
8 hours using 20 million noisy qubits. \emph{Quantum} 5:433, 2021.

\bibitem{goppa1970}
V.~D.~Goppa. A new class of linear correcting codes.
\emph{Problemy Peredachi Informatsii} 6(3):24--30, 1970.

\bibitem{grover1996}
L.~K.~Grover. A fast quantum mechanical algorithm for database
search. En \emph{Proc.\ 28th ACM Symposium on Theory of Computing
(STOC)}, pp.~212--219, 1996.

\bibitem{hemmert2026}
T.~Hemmert. Key recovery for the McEliece cryptosystem using
higher-order vanishing, Cryptology ePrint Archive, Report 2026/1339,
2026.

\bibitem{hemmertwiemers2026}
T.~Hemmert, A.~Wiemers. Distinguishing Goppa codes using
higher-order vanishing. In \emph{Advances in Cryptology ---
CRYPTO 2026}, pp.~191--223, Springer, 2026.

\bibitem{hhk2017}
D.~Hofheinz, K.~H\"ovelmanns, E.~Kiltz. A modular analysis of the
Fujisaki--Okamoto transformation. In \emph{Theory of Cryptography ---
TCC 2017}, LNCS 10677, pp.~341--371, Springer, 2017.

\bibitem{iso18033}
ISO/IEC. \emph{Information technology --- Security techniques ---
Encryption algorithms --- Part~2: Asymmetric ciphers ---
Amendment~2: Classic McEliece}, ISO/IEC 18033-2:2006/Amd~2:2026,
junio de 2026.

\bibitem{kachigar2017}
G.~Kachigar, J.-P.~Tillich. Quantum information set decoding
algorithms. In \emph{Post-Quantum Cryptography --- PQCrypto 2017},
LNCS 10346, pp.~69--89, Springer, 2017.

\bibitem{kirshanova2018}
E.~Kirshanova. Improved quantum information set decoding. In
\emph{Post-Quantum Cryptography --- PQCrypto 2018}, LNCS 10786,
pp.~507--527, Springer, 2018.

\bibitem{kirshanova2022}
E.~Kirshanova, A.~May. Decoding McEliece with a hint --- secret
Goppa key parts reveal everything. In \emph{Security and Cryptography
for Networks --- SCN 2022}, LNCS 13409, Springer, 2022.

\bibitem{lee1988}
P.~J.~Lee, E.~F.~Brickell. An observation on the security of
McEliece's public-key cryptosystem. In \emph{Advances in Cryptology
--- EUROCRYPT~'88}, LNCS 330, pp.~275--280, Springer, 1988.

\bibitem{leon1988}
J.~S.~Leon. A probabilistic algorithm for computing minimum weights
of large error-correcting codes. \emph{IEEE Transactions on
Information Theory} 34(5):1354--1359, 1988.

\bibitem{lidw1994}
Y.~X.~Li, R.~H.~Deng, X.~M.~Wang. On the equivalence of McEliece's
and Niederreiter's public-key cryptosystems. \emph{IEEE Transactions
on Information Theory} 40(1):271--273, 1994.

\bibitem{lidl}
R.~Lidl, H.~Niederreiter. \emph{Finite Fields}, Encyclopedia of
Mathematics and its Applications 20, Cambridge University Press,
2.ª edici\'on, 1997.

\bibitem{msbook}
F.~J.~MacWilliams, N.~J.~A.~Sloane. \emph{The Theory of
Error-Correcting Codes}, North-Holland, 1977.

\bibitem{mmt2011}
A.~May, A.~Meurer, E.~Thomae. Decoding random linear codes in
$\tilde{\mathcal{O}}(2^{0.054n})$. In \emph{Advances in Cryptology ---
ASIACRYPT 2011}, LNCS 7073, pp.~107--124, Springer, 2011.

\bibitem{mo2015}
A.~May, I.~Ozerov. On computing nearest neighbors with applications
to decoding of binary linear codes. In \emph{Advances in Cryptology
--- EUROCRYPT 2015}, LNCS 9056, pp.~203--228, Springer, 2015.

\bibitem{mceliece1978}
R.~J.~McEliece. A public-key cryptosystem based on algebraic coding
theory. \emph{DSN Progress Report 42-44}, Jet Propulsion Laboratory,
pp.~114--116, 1978.

\bibitem{nist8547}
D.~Moody et al. \emph{Transition to post-quantum cryptography
standards}, NIST Internal Report 8547 (borrador p\'ublico inicial),
noviembre de 2024.

\bibitem{niederreiter1986}
H.~Niederreiter. Knapsack-type cryptosystems and algebraic coding
theory. \emph{Problems of Control and Information Theory}
15(2):159--166, 1986.

\bibitem{overbeck2009}
R.~Overbeck, N.~Sendrier. Code-based cryptography. En
\emph{Post-Quantum Cryptography} (D.~J.~Bernstein, J.~Buchmann,
E.~Dahmen, eds.), pp.~95--145, Springer, 2009.

\bibitem{patterson1975}
N.~J.~Patterson. The algebraic decoding of Goppa codes. \emph{IEEE
Transactions on Information Theory} 21(2):203--207, 1975.

\bibitem{persichetti2012}
E.~Persichetti. \emph{Improving the efficiency of code-based
cryptography}, tesis doctoral, University of Auckland, 2012.

\bibitem{prange1962}
E.~Prange. The use of information sets in decoding cyclic codes.
\emph{IRE Transactions on Information Theory} 8(5):5--9, 1962.

\bibitem{randriam2024}
H.~Randriambololona. The syzygy distinguisher. In \emph{Advances in
Cryptology --- EUROCRYPT 2025}, LNCS 15606, Springer, 2025
(revisi\'on ampliada: Cryptology ePrint Archive, Report 2024/1193).

\bibitem{saarinen2026a}
M.-J.~O.~Saarinen. Bit operation cost of ``holdout'' key-recovery
attacks against Classic McEliece, Cryptology ePrint Archive,
Report 2026/1786, 2026.

\bibitem{saarinen2026b}
M.-J.~O.~Saarinen. Two-anchor holdout/Hermite: solving the TII-254
McEliece key recovery challenge, Cryptology ePrint Archive,
Report 2026/1986, 2026.

\bibitem{saarinen2026hover}
M.-J.~O.~Saarinen. HOVER: Higher-order vanishing endomorphism
recovery, Cryptology ePrint Archive, Report 2026/1778, 2026
(NORDSEC 2026, LNCS, en prensa).

\bibitem{sendrier2000}
N.~Sendrier. Finding the permutation between equivalent linear
codes: the support splitting algorithm. \emph{IEEE Transactions on
Information Theory} 46(4):1193--1203, 2000.

\bibitem{shor1994}
P.~W.~Shor. Algorithms for quantum computation: discrete logarithms
and factoring. En \emph{Proc.\ 35th IEEE Symposium on Foundations of
Computer Science (FOCS)}, pp.~124--134, 1994.

\bibitem{shoup2001}
V.~Shoup. A proposal for an ISO standard for public key encryption,
Cryptology ePrint Archive, Report 2001/112, 2001.

\bibitem{stern1989}
J.~Stern. A method for finding codewords of small weight. En
\emph{Coding Theory and Applications}, LNCS 388, pp.~106--113,
Springer, 1989.

\bibitem{vlint}
J.~H.~van Lint. \emph{Introduction to Coding Theory}, Graduate
Texts in Mathematics 86, Springer, 3.ª edici\'on, 1999.

\bibitem{vedenev2026}
K.~Vedenev. Extending distinguishing to key recovery for subfield
subcodes of GRS codes, Cryptology ePrint Archive, Report 2026/1747,
2026.

\bibitem{weis2026}
S.~A.~Weis. Improving GIJS key recovery for Classic McEliece,
Cryptology ePrint Archive, Report 2026/1984, 2026.

\end{thebibliography}
